% 第二十节：洛伦兹变换与帕斯卡定理
\chapter{洛伦兹变换与帕斯卡定理}

\section{二次点列上的射影变换}

本节讨论旋量空间上的一般线性映射 $\varphi$。它表示二次点列上的一个射影变换，代数上由一个
旋量张量 ${t^\alpha}_\beta$ 表示，其特点是 $t_{\alpha\beta}$ 并不对称。对称的情形对应一个
对合，已在二次点列与二次对合一节中讨论过；非对称的情形不能再写成矢量的形式，只有借助旋量
才能处理。

设 $\mathfrak a$ 是 $S$ 中的旋量，$\mathfrak a' =\varphi(\mathfrak a)$ 由
\begin{equation}\label{eq:varphi}
\mathfrak a'{}^\alpha=t^\alpha{}_\beta\,\mathfrak a^\beta
\end{equation}
给出。约定用 $a'$ 表示 $\varphi(a)$，即若 $a=\mathfrak a\mathfrak a$，则
$a'=\mathfrak a'\mathfrak a'$。由于 $g$ 上的点与射影旋量一一对应，$\varphi$ 便是二次点列上的
一个射影变换；$\varphi$ 的不动点对应于 $t^\alpha{}_\beta$ 的特征旋量。

\section{命题}

\begin{proposition}\label{prop:lorentz}
$\varphi$ 可以从二次点列上的射影变换线性延拓为全平面上的射影变换，这一射影变换保持 $g$ 以及
点乘不变。反过来，每一个保持圆锥曲线（的像）不变的射影变换都可以用这样的 $\varphi$ 表示。
出于它与物理的相似性，称它为\emph{洛伦兹变换}。
\end{proposition}

\begin{remark}[与相对论速度空间的类比]
命题 \ref{prop:lorentz} 的一个应用是在相对论粒子碰撞中。把射影平面看作四维动量商掉模长后
得到的粒子速度空间（商掉静质量），则 $\varphi$ 代表所有可能的参考系变换：$g$ 上的点代表
光子，$g$ 内部的点代表有静质量的粒子。于是对一个粒子碰撞问题，可以画出它的速度图，再用
后面的命题 2、3、4 处理。
\end{remark}

\begin{proposition}\label{prop:locus}
设 $a$ 是 $g$ 上的动点，则直线 $a\times a'$ 的极点 $K$ 的轨迹构成一条圆锥曲线 $h$，
$h$ 与 $g$ 必定相切于两点 $p,q$。反之，每给定一条与 $g$ 相切于两点的圆锥曲线 $h$，都可以
通过这种作图方法构造出对应的线性映射 $\varphi$ 及其逆 $\varphi^{-1}$。
\end{proposition}

映射 $\varphi$ 与洛伦兹变换比较抽象，不容易看出几何含义；借助命题 \ref{prop:locus}，
可以用一类特殊的圆锥曲线来构造 $\varphi$，从而把它变得直观。

\begin{proposition}\label{prop:construction-outside}
设 $a$ 是 $g$ 外一点，则可以通过如下步骤得到 $a'$：
\begin{enumerate}[label=(\arabic*),leftmargin=2.4em]
\item 作 $a$ 到 $g$ 的两条切线，切点分别为 $b,c$；
\item $\varphi$ 在 $g$ 上有定义，故可作出 $b',c'$；
\item 作 $b',c'$ 处的切线，其交点即为 $a'$。
\end{enumerate}
\end{proposition}

命题 \ref{prop:construction-outside} 给出了 $\varphi$ 延拓的一种作图方式。

\begin{proposition}\label{prop:family}
用 $C$ 表示所有与 $g$ 相切于 $p,q$ 两点的圆锥曲线构成的集合。对任意一点 $a$，在 $C$ 中
存在唯一一条经过 $a$ 的圆锥曲线，记作 $c(a)$，并且恒有
\[
c(a)=c(a') .
\]
换言之，映射 $\varphi$ 把任意一条与 $g$ 相切于 $p,q$ 两点的圆锥曲线映为自身。于是所有
这样的圆锥曲线在平面上划出一族 $\varphi$ 轨道。
\end{proposition}

\begin{proposition}\label{prop:construction-inside}
设 $a$ 不在 $g$ 上，则可以通过如下步骤作出 $a'$：
\begin{enumerate}[label=(\arabic*),leftmargin=2.4em]
\item 作直线 $a\times p$，交 $g$ 于 $b$；
\item 作出 $b'$；
\item 作直线 $b'\times p$，它与 $c(a)$ 的交点即为 $a'$。
\end{enumerate}
\end{proposition}

\begin{proposition}[帕斯卡定理]\label{prop:pascal}
\begin{enumerate}[label=(\arabic*),leftmargin=2.4em]
\item $g$ 上的三对对应点 $a,a',b,b',c,c'$ 能唯一确定射影变换 $\varphi$，使得
      $a'=\varphi(a)$，$b'=\varphi(b)$，$c'=\varphi(c)$。
\item 设 $s=p\times q$，称 $s$ 为 $\varphi$ 对应的\emph{帕斯卡线}。设 $a,b$ 是 $g$ 上的
      任意两点，则点 $(a\times b')\times(a'\times b)$ 必定在帕斯卡线上。
\end{enumerate}
\end{proposition}

\section{命题的证明}

\subsection{命题 \ref{prop:lorentz} 的证明}

根据 \eqref{eq:varphi}，$\mathfrak a'{}^\alpha=t^\alpha{}_\beta\mathfrak a^\beta$。设
$a=\mathfrak a\mathfrak a$，则由旋量与类光矢量的对应关系有 $a'=\mathfrak a'\mathfrak a'$，
即
\[
a'^{\alpha\dot\alpha}
=t^\alpha{}_\beta\mathfrak a^\beta\,
t^{\dot\alpha}{}_{\dot\beta}\mathfrak a^{\dot\beta}
=t^\alpha{}_\beta t^{\dot\alpha}{}_{\dot\beta}a^{\beta\dot\beta},
\]
这是关于 $a^{\beta\dot\beta}$ 的线性映射。此时 $a^{\beta\dot\beta}$ 作为定义域只包含
$\det a^{\beta\dot\beta}=0$ 的部分，但线性映射
$t^\alpha{}_\beta t^{\dot\alpha}{}_{\dot\beta}a^{\beta\dot\beta}$ 可自然地把定义域延拓到所有
对称旋量张量，因而延拓到平面上的所有点。用矢量指标写出即
\[
a'^i=\sigma^i_{\alpha\dot\alpha}t^\alpha{}_\beta t^{\dot\alpha}{}_{\dot\beta}
\sigma_j^{\beta\dot\beta}a^j,
\]
可见它确实是全平面上的射影变换。

又因为 $\varphi$ 把 $g$ 上的点映为 $g$ 上的点，所以它保持 $g$ 的像（圆锥曲线对应的点集）
不变；圆锥曲线张量由其像完全确定，因此 $g$ 作为张量也关于 $\varphi$ 不变。于是由 $g_{ij}$
定义的点乘也关于 $\varphi$ 不变。反过来，在二维空间加一维时间中保持时空距离不变的参考系
变换是三维的（两个 boost 与一个空间旋转），而二维矢量空间中的射影变换也是三维的，
两者恰好一一对应。

\subsection{命题 \ref{prop:locus} 的证明}

只需用 $t_{\alpha\beta}$ 表示出 $h_{\alpha\beta\gamma\delta}$，即找出这样的
$h_{\alpha\beta\gamma\delta}$，使得对任意的 $\mathfrak c$ 恒有
\begin{equation}\label{eq:hlocus}
h_{\alpha\beta\gamma\delta}\,
t^\beta{}_{\dot\beta}t^\delta{}_{\dot\delta}\,
\mathfrak c^\alpha\mathfrak c^{\dot\beta}\mathfrak c^\gamma\mathfrak c^{\dot\delta}=0 .
\end{equation}

首先注意，射影变换 $t^\alpha{}_\beta$ 的不动点方程是二次的，因此必定有两个不动点，记为
$\mathfrak a,\mathfrak b$，即
\[
\mathfrak a^\alpha=t^\alpha{}_\beta\mathfrak a^\beta,\qquad
\mathfrak b^\alpha=t^\alpha{}_\beta\mathfrak b^\beta .
\]
先看第一个方程。把它降指标得 $\mathfrak a_\alpha=t_{\alpha\beta}\mathfrak a^\beta$，
两边乘以 $\mathfrak a^\alpha$，利用 $\mathfrak a^\alpha\mathfrak a_\alpha=0$，得到
$t_{\alpha\beta}\mathfrak a^\alpha\mathfrak a^\beta=0$。设
\[
s_{\alpha\beta}=t_{\alpha\beta}+t_{\beta\alpha},
\]
则 $s_{\alpha\beta}\mathfrak a^\alpha\mathfrak a^\beta=0$；同理
$s_{\alpha\beta}\mathfrak b^\alpha\mathfrak b^\beta=0$。由于 $s_{\alpha\beta}$ 对称，它表示
一个对合变换，而 $\mathfrak a,\mathfrak b$ 是其不动点，故
\begin{equation}\label{eq:s}
s_{\alpha\beta}=t_{\alpha\beta}+t_{\beta\alpha}=(\mathfrak a\mathfrak b)_{\alpha\beta},
\end{equation}
它对应直线 $p\times q$。

因为 $g,h$ 相切于两点，所以 $g,h,ss$ 构成二次线系，于是可设
\begin{equation}\label{eq:hansatz}
h_{\alpha\beta\gamma\delta}
=s_{\alpha\beta}s_{\gamma\delta}+\lambda g_{\alpha\beta\gamma\delta},
\end{equation}
其中 $\lambda$ 待定。把 \eqref{eq:hansatz} 代入 \eqref{eq:hlocus}，得
\[
\bigl(s_{\alpha\beta}s_{\gamma\delta}t^\beta{}_{\dot\beta}t^\delta{}_{\dot\delta}
+\lambda g_{\alpha\beta\gamma\delta}t^\beta{}_{\dot\beta}t^\delta{}_{\dot\delta}\bigr)
\mathfrak c^\alpha\mathfrak c^{\dot\beta}\mathfrak c^\gamma\mathfrak c^{\dot\delta}=0 .
\]
小括号中第一项为
\[
s_{\alpha\beta}s_{\gamma\delta}t^\beta{}_{\dot\beta}t^\delta{}_{\dot\delta}
=(t_{\alpha\beta}+t_{\beta\alpha})(t_{\gamma\delta}+t_{\delta\gamma})
t^{\beta\beta'}t^{\delta\delta'}e_{\beta'\dot\beta}e_{\delta'\dot\delta}
=\bigl({{t^2}_\alpha}^{\beta'}+\delta_\alpha^{\beta'}\bigr)
\bigl({{t^2}_\gamma}^{\delta'}+\delta_\gamma^{\delta'}\bigr)
e_{\beta'\dot\beta}e_{\delta'\dot\delta},
\]
即
\[
s_{\alpha\beta}s_{\gamma\delta}t^\beta{}_{\dot\beta}t^\delta{}_{\dot\delta}
=\bigl({t^2}_{\alpha\dot\beta}+e_{\alpha\dot\beta}\bigr)
\bigl({t^2}_{\gamma\dot\delta}+e_{\gamma\dot\delta}\bigr),
\qquad
{{t^2}{}^\alpha}_\gamma={t^\alpha}_\beta{t^\beta}_\gamma .
\]
把它乘以 $\mathfrak c^\alpha\mathfrak c^{\dot\beta}\mathfrak c^\gamma\mathfrak c^{\dot\delta}$：
含有 $e_{\alpha\dot\beta}$ 的项因 $e_{\alpha\dot\beta}\mathfrak c^\alpha\mathfrak c^{\dot\beta}=0$
而消失，留下的只有
\[
{t^2}_{\alpha\dot\beta}{t^2}_{\gamma\dot\delta}
\mathfrak c^\alpha\mathfrak c^{\dot\beta}\mathfrak c^\gamma\mathfrak c^{\dot\delta}
=\bigl({t^2}_{\alpha\dot\beta}\mathfrak c^\alpha\mathfrak c^{\dot\beta}\bigr)^2 .
\]
小括号中的第二项为
\[
g_{\alpha\beta\gamma\delta}t^\beta{}_{\dot\beta}t^\delta{}_{\dot\delta}
=(e_{\alpha\gamma}e_{\beta\delta}+e_{\alpha\delta}e_{\beta\gamma})
t^\beta{}_{\dot\beta}t^\delta{}_{\dot\delta}
=e_{\alpha\gamma}e_{\dot\delta\dot\beta}-t_{\alpha\dot\delta}t_{\gamma\dot\beta},
\]
其中第一项乘以 $\mathfrak c^\alpha\mathfrak c^{\dot\beta}\mathfrak c^\gamma\mathfrak c^{\dot\delta}$
后为零，故只留下
\[
(-t_{\alpha\dot\delta}t_{\gamma\dot\beta})\,
\mathfrak c^\alpha\mathfrak c^{\dot\beta}\mathfrak c^\gamma\mathfrak c^{\dot\delta}
=-\bigl(t_{\alpha\dot\beta}\mathfrak c^\alpha\mathfrak c^{\dot\beta}\bigr)^2 .
\]
于是 \eqref{eq:hlocus} 等价于
\begin{equation}\label{eq:hcond}
\bigl({t^2}_{\alpha\beta}\mathfrak c^\alpha\mathfrak c^\beta\bigr)^2
-\lambda\bigl(t_{\alpha\beta}\mathfrak c^\alpha\mathfrak c^\beta\bigr)^2=0 .
\end{equation}

表面上看，无论 $\lambda$ 如何取值也不能使 \eqref{eq:hcond} 成为恒等式，但事实并非如此。
$2\times2$ 矩阵 $A$ 满足恒等式
\begin{equation}\label{eq:cayley}
A\tr A-A^2=(\det A)I,
\end{equation}
即 $A^2-(\tr A)A+(\det A)I=0$。把 \eqref{eq:cayley} 应用于 ${t^\alpha}_\beta$，得到
\[
{t^\alpha}_\beta{t^\mu}_\mu-{{t^2}{}^\alpha}_\beta=(\det t)\delta^\alpha_\beta,
\]
也就是
\[
t_{\alpha\beta}{t^\mu}_\mu-{t^2}_{\alpha\beta}=(\det t)e_{\alpha\beta}.
\]
两边同乘 $\mathfrak c^\alpha\mathfrak c^\beta$，注意 $e_{\alpha\beta}\mathfrak c^\alpha\mathfrak c^\beta=0$，
便得
\[
{t^\mu}_\mu\bigl(t_{\alpha\beta}\mathfrak c^\alpha\mathfrak c^\beta\bigr)
={t^2}_{\alpha\beta}\mathfrak c^\alpha\mathfrak c^\beta .
\]
代入 \eqref{eq:hcond} 可知，只需取
\[
\lambda=({t^\mu}_\mu)^2,
\]
\eqref{eq:hcond} 便成为恒等式，从而证明了 $a\times a'$ 的极点确实在圆锥曲线 $h$ 上。于是
\begin{equation}\label{eq:hformula}
\boxed{\;h_{\alpha\beta\gamma\delta}
=(t_{\alpha\beta}+t_{\beta\alpha})(t_{\gamma\delta}+t_{\delta\gamma})
+({t^\mu}_\mu)^2g_{\alpha\beta\gamma\delta}.\;}
\end{equation}

最后注意，圆锥曲线共有 $5$ 个自由度，两个相切条件产生 $2$ 个约束，因此这样的 $h$ 共有
$3$ 个自由度，等于射影变换的自由度个数，故 $h$ 与 $\varphi$ 相互对应。

\subsection{命题 \ref{prop:construction-outside} 的证明}

利用 $\varphi$ 保持点乘这一性质。由 $a\cdot b=0$ 得 $a'\cdot b'=0$，同理
$a'\cdot c'=0$。因此 $a'$ 同时在 $b'$ 与 $c'$ 的极线上，即
\[
a'=b'\times c',
\]
这正是 $b',c'$ 处两条切线的交点。

\subsection{命题 \ref{prop:family} 的证明}

先证唯一性。五点确定一条圆锥曲线；"相切"可以看作同时经过切点附近切线上无限接近的两个点，
于是命题中的问题可以转化为五点确定一条圆锥曲线的问题，因此 $c(a)$ 唯一。

再证 $c(a)=c(a')$。由命题 \ref{prop:locus}，$C$ 中的元素都可用
\begin{equation}\label{eq:family}
c_\mu=s_{\alpha\beta}s_{\gamma\delta}+\mu\,g_{\alpha\beta\gamma\delta},
\qquad \mu\in\R,
\end{equation}
描述，其中 $s$ 由 \eqref{eq:s} 给出，对应直线 $p\times q$。由于 $p,q$ 是 $\varphi$ 的不动
点，$g$ 与 $s$ 在 $\varphi$ 下分别变为自身的倍数，因此 $c_\mu$ 的像仍属于曲线系
$\{ss+\mu'g\}$，即
\[
c_\mu'=\alpha\,s_{\alpha\beta}s_{\gamma\delta}
+\beta\mu\,g_{\alpha\beta\gamma\delta},
\]
其中 $\alpha,\beta$ 是与 $\mu$ 无关的常数。

取 $\mu=\lambda=({t^\mu}_\mu)^2$，此时 $c_\lambda$ 正是命题 \ref{prop:locus} 中的圆锥曲线
$h$。由 \eqref{eq:hformula} 与命题 \ref{prop:locus}，$h$ 的像就是它自身（作为点集），于是
\[
\alpha\,ss+\beta\lambda\,g
\ \text{与}\ ss+\lambda\,g\ \text{成正比}.
\]
比较系数（$\lambda\neq0$ 时）得 $\alpha=\beta$。代入 \eqref{eq:family} 的一般 $\mu$，
\[
c_\mu'=\alpha\,ss+\alpha\mu\,g=\alpha\,(ss+\mu g)=\alpha\,c_\mu,
\]
即 $c_\mu'$ 与 $c_\mu$ 表示同一条圆锥曲线，因而 $c(a')=c(a)$。当 $\lambda=0$ 时，
$\operatorname{tr}t=0$，$h$ 退化为 $ss$；此时结论可由 $\operatorname{tr}t\neq0$ 的情形取
极限得到。

\begin{remark}
$C$ 中的元素有两条是退化的：一条是 $ss$，另一条是 $p,q$ 处两条切线的对称积。设后者两条
切线的交点为 $r$，则 $r$ 是 $s$ 的极点。在整个平面上，$\varphi$ 的不动点有且仅有 $p,q,r$
三个。
\end{remark}

\subsection{命题 \ref{prop:construction-inside} 的证明}

只需证明按题中三步作图后，$a'$ 位于直线 $b'\times p$ 上。由作图，$b$ 在直线 $a\times p$ 上，
故
\[
(b\times p)\cdot a=0 .
\]
射影变换保持 Levi-Civita 记号（因而保持叉乘）与点乘，于是上式经 $\varphi$ 作用后成为
\[
(b'\times p')\cdot a'=0 .
\]
又 $p'=p$，所以 $(b'\times p)\cdot a'=0$，即 $a'$ 在直线 $b'\times p$ 上。结合命题
\ref{prop:family} 给出的 $a'\in c(a)$，二者的交点即为 $a'$。

\begin{remark}
$p$ 只是两个不动点中的一个。对另一个不动点 $q$ 同样操作：作直线 $a\times q$ 交 $g$ 于
$c$，作出 $c'$，再作直线 $c'\times q$，同样可证 $(c'\times q)\cdot a'=0$。因此
\[
a'=(b'\times p)\times(c'\times q).
\]
\end{remark}

\subsection{命题 \ref{prop:pascal} 的证明}

(1) 令 $g$ 上任意点 $d$ 的像 $d'$ 满足 $a,b,c,d$ 的交比等于 $a',b',c',d'$ 的交比，
这便唯一地确定了映射 $\varphi$。

(2) 由二次点列与二次对合中的结论，三条直线 $a\times b'$、$a'\times b$ 与 $s=p\times q$
共点，当且仅当存在对合 $f$，使得 $(p,q)$、$(a,b')$、$(a',b)$ 均关于 $f$ 对合。

用旋量语言写出：$a=\mathfrak a\mathfrak a$，$a'=\mathfrak a'\mathfrak a'$，
$b=\mathfrak b\mathfrak b$，$b'=\mathfrak b'\mathfrak b'$，于是
\[
a\times b'=(\mathfrak a\mathfrak b'),\qquad a'\times b=(\mathfrak a'\mathfrak b).
\]
在本节命题 \ref{prop:locus} 的证明中已得到
$s_{\alpha\beta}=t_{\alpha\beta}+t_{\beta\alpha}$。对合 $f$ 的三个条件分别为
\begin{equation}\label{eq:pascaleq}
f_{\alpha\beta}\mathfrak a^\alpha\mathfrak b'{}^\beta=0,\qquad
f_{\alpha\beta}\mathfrak a'{}^\alpha\mathfrak b^\beta=0,\qquad
f_{\alpha\beta}t^{\alpha\beta}=0 .
\end{equation}
前两个方程利用 $\mathfrak a'{}^\alpha=t^\alpha{}_\beta\mathfrak a^\beta$、
$\mathfrak b'{}^\beta=t^\beta{}_\gamma\mathfrak b^\gamma$ 展开为
\[
f_{\alpha\beta}\mathfrak a^\alpha t^\beta{}_{\dot\beta}\mathfrak b^{\dot\beta}=0,
\qquad
f_{\alpha\beta}t^\alpha{}_{\dot\alpha}\mathfrak a^{\dot\alpha}\mathfrak b^\beta=0 .
\]

为证明这三个方程可以同时成立，先令第三个方程成立，再证明此时前两个方程等价。
改设
\[
h_\alpha{}^\gamma=f_{\alpha\beta}t^{\beta\gamma},
\]
则第三个方程 $f_{\alpha\beta}t^{\alpha\beta}=0$ 改写为 $h_\gamma{}^\gamma=0$；前两个方程分别
改写为
\[
h_{\alpha\dot\beta}\mathfrak a^\alpha\mathfrak b^{\dot\beta}=0,
\qquad
h_{\beta\dot\alpha}\mathfrak a^{\dot\alpha}\mathfrak b^\beta=0
\]
（改写第二个方程时利用了 $f_{\alpha\beta}=f_{\beta\alpha}$）。而由 $h_\gamma{}^\gamma=0$
可以推出 $h_{\alpha\beta}=h_{\beta\alpha}$：这是二维矩阵的结论，零迹决定了对称。于是前两个
方程完全等价，只剩下一个方程。对称无迹的 $h$ 有两个独立分量，一个线性方程将其确定到相差
一个常数因子，因而对合 $f$ 被唯一确定。证毕。
